\documentclass[11pt]{article}

\usepackage[margin=1in]{geometry}
\usepackage{amsmath}
\usepackage{booktabs}
\usepackage{enumitem}
\usepackage{graphicx}
\usepackage{siunitx}
\usepackage{hyperref}
\usepackage{gensymb}


\title{Advanced Rolling Test}
\author{Rift / Team 09}
\date{\today}

\begin{document}


\maketitle

\section*{Test ID}
\textbf{T-01AR Advanced Rolling Test} 

\section*{Test Objective}
Determine the limits of possible rolling techniques, including lateral rotation (cartwheeling), or successfully performing a series of different rolling sequences (somersaulting, or typical direct rolling, into other forms of rolling).

\section*{Robot Configuration}
\begin{itemize}[noitemsep]
    \item Robot configurations: Standing
    \item Tube pressures: 15 psi
    \item Control mode: Open-loop
    \item Electronics status: Batteries are fully charged
\end{itemize}

\section*{Materials}
\begin{itemize}[noitemsep]
    \item None
    
\end{itemize}

\section*{Environment and Setup}
\begin{itemize}[noitemsep]
    \item Surface type: Finished Smooth Concrete
    \item Surface angle: $\psi = 0 {\degree}$ 
    \item External load: None 
    \item Measurement method: Noting differentiation between the expected sequence of rolling motion in simulation compared to what the rover does.
\end{itemize}

\section*{Initial Conditions}
\begin{itemize}[noitemsep]
    \item Initial position: $x_0 = 0$
    \item Initial orientation: $\theta_z = 0 {\degree}$
    \item Initial shape state: Standing
\end{itemize}

\section*{Command Sequence}
\begin{itemize}[noitemsep]
    \item Command type: \underline{\hspace{5 cm}}
    \item Command values: \underline{\hspace{5 cm}}
    \item Update rate: 3~Hz
    \item Command Steps: 8~s
\end{itemize}

\section*{Measured Quantities}
\begin{itemize}[noitemsep]
    \item None
\end{itemize}

\section*{Success Criteria (Per Trial)}
\begin{itemize}[noitemsep]
    \item Rover passes or fails on the sequences that are sent to cause it to roll.
\end{itemize}

\section*{Test Procedure}
\begin{enumerate}[noitemsep]
    \item Place robot in standing position
    \item Complete pretest checklist
    \item Send command to rover to perform sequence to roll in different combinations.

\end{enumerate}

\section*{Number of Trials}
\begin{itemize}[noitemsep]
    \item Total trials: 1
    \item Time between trials: N/A
\end{itemize}

\section*{Results}
Summarize the observed performance metrics.

\begin{center}
\begin{tabular}{lc}
\toprule
Roll Sequence& Pass/Fail\\
 &\\
\midrule
&  \\
\end{tabular}
\end{center}

\section*{Observed Failure Modes}
Describe any observed failure behaviors.

\begin{itemize}[noitemsep]
    \item Truss buckles during motion
    \item Rover is not capable of rolling
    \item Rover positioning is notably distinct from the expected position in simulation.
\end{itemize}

\section*{Conclusions}
Provide a concise interpretation of the results.

\textit{Example:} The robot achieved consistent forward motion on flat ground with an average speed of \underline{\hspace{1 cm}}~cm/s. Performance degraded beyond a slope of \underline{\hspace{1 cm}}~\si{\degree}, where sustained forward progress was not achieved.

\section*{Notes and Limitations}
Document sources of variability or limitations.

\end{document}